\section{Introduction}\label{sec:introduction}

For a complete connected smooth Riemannian manifold $(M^n,g)$ without
boundary, let $\mathcal H_d(M,g)$ be the real vector space of harmonic
functions satisfying
\[
 |u(x)|\leq C_u(1+d_g(o,x))^d\qquad(x\in M)
\]
for some finite constant $C_u$, where $d\geq0$ and $o\in M$ is a base
point. Changing $o$ does not change the space. Write
$h_d(M,g)=\dim_{\mathbb R}\mathcal H_d(M,g)$. On Euclidean space, an
entire harmonic function with integer growth at most $k$ is a harmonic
polynomial of degree at most $k$. Consequently
\begin{equation}\label{eq:euclidean-comparator}
 h_k(\mathbb R^n,g_{\mathrm E})
 =\binom{n+k-1}{k}+\binom{n+k-2}{k-1}\qquad(k\geq1).
\end{equation}
The integer-degree comparison problem asks whether $\operatorname{Ric}_g
\geq0$ forces $h_k(M,g)\leq h_k(\mathbb R^n,g_{\mathrm E})$ for every
integer $k\geq1$. The condition defining $\mathcal H_k$ is pointwise at
every distance. Bounds only along a sequence of spheres do not establish it.

\begin{theorem}\label{thm:main}
There are an even integer $n\geq8$, an integer $k\geq2$, and a complete
smooth metric $g$ on $\mathbb R^n$ such that $\operatorname{Ric}_g\geq0$
and
\[
 h_k(\mathbb R^n,g)>h_k(\mathbb R^n,g_{\mathrm E}).
\]
The metric is Euclidean near the origin, and its asymptotic volume ratio
satisfies
\[
 0<\lim_{R\to\infty}
 \frac{\operatorname{vol}_g B_g(0,R)}{\omega_nR^n}<1,
\]
where $\omega_n$ is the volume of the Euclidean unit ball.
\end{theorem}

The exact finite spectral selection in Appendix~\ref{sec:finite-selection}
permits the choices \(n=16\) and \(k=50000\).

Theorem~\ref{thm:main} answers the integer-degree form of Yau's dimension
comparison question in the negative. Yau's questions on polynomial-growth harmonic functions distinguish finite
dimensionality from the sharper Euclidean comparison~\cite{Yau,LiTam}.
Li and Tam proved the sharp linear-growth bound $h_1(M,g)\leq n+1$ under
nonnegative Ricci curvature~\cite{LiTam}. Colding and Minicozzi proved
finite dimensionality at every fixed degree and later obtained bounds of
the optimal order $d^{n-1}$ as $d$ grows~\cite{CM97,CM98}. Those estimates
leave open the exact Euclidean count at a particular integer degree.

Donnelly constructed counterexamples to the analogous comparison at a real
degree strictly between one and two in dimensions at least
five~\cite{Donnelly,CaiLai}. That fact alone does not decide the integer
endpoint two: the Euclidean space of quadratic harmonic polynomials is
larger than its space at a degree between one and two. Cai and Lai make
this distinction explicit and establish sharp comparison with the
additional hypothesis of local conformal flatness~\cite{CaiLai}.
The theorem above addresses the integer question under the Ricci
hypothesis alone.

Cone-spectral results explain another boundary of the problem. For complete
manifolds with nonnegative Ricci curvature, maximal volume growth, and a
unique tangent cone at infinity, Huang's
Theorem~1.6 relates the harmonic dimension to the spectrum of the cone
link~\cite{Huang}. Our metric instead has round and nonround subsequential
cone limits. Nonunique cones with positive Ricci curvature and Euclidean
volume growth already occur in Perelman's construction~\cite{Perelman};
Colding and Naber developed a method for slowly varying link
families with a common volume form~\cite{ColdingNaber}. We use the same geometric possibility of a
moving link, while proving the curvature estimates for our particular
deformations. Xu's nonunique-cone three-circles theorem requires a comparison
exponent outside the union of cone-degree spectra~\cite[Theorem~3.4]{Xu}.
A separate geometric conclusion of our construction is that this union
contains a neighborhood of the selected integer $k$, so the theorem cannot
be applied at $k$ or at sufficiently nearby exponents.

\subsection*{Why a moving link can change the count}

Put $n=m+1$. On the metric cone $dr^2+r^2h$ over an $m$-dimensional link,
an angular eigenfunction of the nonnegative Laplacian with eigenvalue
$\lambda$ has a growing homogeneous harmonic mode $r^d\phi$, where
\[
 d(d+m-1)=\lambda,\qquad d\geq0.
\]
For one fixed link, each angular direction therefore carries one fixed
growth rate. Our construction lets a harmonic function use several rates
in succession. The challenge is to arrange that succession exactly:
an arbitrarily small component in an unwanted faster direction can dominate
after a sufficiently long interval.

The links are volume-normalized Berger metrics on odd spheres
$S^m$, $m=2s-1$. Their distinguished Hopf direction is allowed to change.
Differentiation along this circle direction is the Hopf generator; on
complexified harmonics its eigenvalues are \(\mathrm i\) times the
integer Hopf charges. Every angular Laplacian nevertheless preserves the same finite-dimensional
space $V_l$ of round spherical harmonics of degree $l$; write
$N_l=\dim_{\mathbb R}V_l$. At a fixed nonround link, write its growing rates
on $V_l$ in increasing order as $d_{l,1},\ldots,d_{l,N_l}$. The spectral
calculation supplies a finite $L\geq2$ and integers $0\leq M_l\leq N_l$
for $2\leq l\leq L$ such that
\begin{equation}\label{eq:overview-selection}
 \sum_{l=2}^L M_l>h_k(\mathbb R^{m+1},g_{\mathrm E}),
 \qquad
 \frac1{M_l}\sum_{i=1}^{M_l}d_{l,i}<k\quad(M_l>0).
\end{equation}
Thus the count will be large enough if each selected harmonic function can
spend asymptotically equal amounts of logarithmic radius at all the selected
rates in its block. The inequality for the mean is strict; individual rates
may exceed $k$.

The source of \eqref{eq:overview-selection} is the distribution of Hopf
charges inside $V_l$. The squared charges divided by $l^2$ tend to a
beta distribution. The mean of the lowest fraction of rates is described by
an integral of its decreasing quantile. In a sufficiently large odd link
dimension, a weighted quantile inequality shows that selecting directions by
their mean produces a strict surplus over the Euclidean count. This fixes the
link parameters, the integer $k$, and a finite degree range before the
radial construction begins. An exact-integer computation gives a separate
finite selection at specific parameters; its role is to verify that finite
spectral input, rather than the subsequent geometric or transmission proof.

\subsection*{From angular control to entire harmonic functions}

Short deformations near a round link provide the exchanges of directions.
For every fixed finite range $2\leq l\leq L$, the trace-free squares of
the Hopf generators, using all orthogonal complex structures, generate
\[
 \bigoplus_{l=2}^L\mathfrak{sl}(V_l).
\]
The real form and the direct sum matter: one must choose the same finite
sequence of geometric pulses while independently arranging a prescribed
signed permutation in each $V_l$. A finite product of exponentials realizes
these permutations at a parameter value where the word map has invertible
differential. The transmission argument below shows that this algebraic
freedom survives the harmonic equation.

Write the metric as $g=dr^2+r^2\gamma(t)$ with $t=\log r$. A harmonic
function in $V_l$ is governed by a second-order matrix equation. Smoothness
at the Euclidean center selects a distinguished regular solution space.
Its logarithmic derivative $P_l$ satisfies a matrix Riccati equation;
positive barriers keep $P_l$ bounded and symmetric. A long round rest makes
$P_l$ approach a scalar matrix, independently of all earlier admissible
pulses. During the next control interval, which is short relative to the
following dwell, the determinant-normalized value transfer converges in
$C^1$ to the finite algebraic word as the period index tends to infinity.

The pulse-end approximation is only an intermediate fact: the outgoing
radial derivative can still inject a small component into a faster direction
during the following long leg, on which the angular eigenvectors remain
fixed. We factor out the positive diagonal transfer \(D_{l,j}\) obtained by entering that entire leg with the round
growing derivative. The remaining multiplier of the derivative error is
bounded independently of the leg's length and anisotropy
(Lemma~\ref{lem:leg}). Thus the full-period map, with this diagonal
factor removed and its determinant normalized, is \(C^1\)-close to the ideal
word, uniformly over all preceding controls. Its local inverse tunes every complete period exactly to
\(s_{l,j}D_{l,j}\Pi_l\), where \(\Pi_l\) cycles the selected axes and
\(s_{l,j}>0\) is a scalar.

The geometry and the growth argument use different features of the same
schedule. In period $j$, the control intervals have length comparable to
$j^6$ and the nonround dwell has length $j^7$, while the accumulated
logarithmic radius is comparable to $j^8$. Angular motion is consequently
slow enough for a small concave radial tail to preserve nonnegative Ricci
curvature. The diagonal leg contributes
$\log(D_{l,j})_{ii}=j^7(d_{l,i}+o(1))$, while
$\log s_{l,j}=O(j^6)$. Exact cycling and the asymptotic equality of finitely
many consecutive dwell lengths yield the means in
\eqref{eq:overview-selection}. Bounded Riccati matrices control the intervals
between period endpoints. Since a whole period is negligible compared with
the accumulated logarithmic radius, the same exponent bound holds at every
radius. The strict inequality below $k$ then gives the pointwise growth
condition, and invertibility of the regular fundamental matrices gives the
claimed number of independent smooth harmonic functions.

All spaces and functions are real unless explicitly complexified. Angular
Laplacians are nonnegative. Matrices act on column vectors, and positivity of
a matrix means positivity as a symmetric operator. Constants may depend on
the fixed finite degree range and fixed compact family of controls; they are
uniform in the period index and in the preceding admissible controls.


Sections~\ref{sec:links} and \ref{sec:counting} establish the Berger formulas
and the finite spectral surplus. Section~\ref{sec:control} proves simultaneous
real control, and Section~\ref{sec:geometry} realizes every admitted control
history by a complete metric with nonnegative Ricci curvature.
Section~\ref{sec:transmission} tunes the full-period transmissions exactly;
Section~\ref{sec:conclusion} proves the every-point growth estimate and the
cone conclusions. Appendix~\ref{sec:finite-selection} records the independent
finite arithmetic selection.
